\documentclass[11pt,a4paper]{article}
\usepackage[utf8]{inputenc}
\usepackage{amsmath,amssymb,amsfonts}
\usepackage{geometry}
\usepackage{hyperref}

\geometry{margin=1in}

\title{\textbf{\Huge Non-Linear Tube Saturation \& Overdrive DSP}\\[0.5em]
\Large Mathematical Models of \texttt{[drive\textasciitilde{}]} and Waveshaping Objects}
\author{\textbf{Time Dilation DAW (v0.0.1) Engineering Group}}
\date{\today}

\begin{document}

\maketitle

\begin{abstract}
This paper formulates the non-linear transfer functions, hyperbolic tangent $\tanh(x)$ waveshaping, asymmetric grid bias tube saturation equations, cubic soft-clipping polynomials, and harmonic gain normalization used in the \texttt{[drive\textasciitilde{}]} node object.
\end{abstract}

\section{Symmetric Hyperbolic Tangent ($\tanh$) Saturation}
In vacuum tube triode amplification, soft saturation limits signal amplitude as input voltage increases. The symmetric transfer function $f(x)$ for input signal $x(n)$ with drive gain $G \ge 1.0$ is:
\begin{equation}
y(n) = \tanh\left( G \cdot x(n) \right)
\end{equation}
Taking the Taylor series expansion of $\tanh(u)$ around $u = 0$:
\begin{equation}
\tanh(u) = u - \frac{u^3}{3} + \frac{2u^5}{15} - \frac{17u^7}{315} + O(u^9)
\end{equation}
The presence of odd-powered polynomial terms ($u^3, u^5, u^7$) introduces odd harmonic overtones ($3f, 5f, 7f$) into the output spectrum, generating warm analog saturation.

\section{Asymmetric Triode Vacuum Tube Bias Modeling}
To simulate authentic vacuum tube asymmetric clipping (which introduces even harmonics $2f, 4f, 6f$ associated with tubey warmth), a DC grid bias offset $V_{\text{bias}} \in [-0.5, 0.5]$ is injected prior to non-linear transformation:
\begin{equation}
y_{\text{raw}}(n) = \tanh\left( G \cdot x(n) + V_{\text{bias}} \right) - \tanh(V_{\text{bias}})
\end{equation}
Subtracting $\tanh(V_{\text{bias}})$ removes the static DC offset, preserving zero DC operating bias at rest.

\section{Cubic Soft-Clipping Polynomial}
For lightweight processing without transcendental functions:
\begin{equation}
y_{\text{cubic}}(x) =
\begin{cases}
-1.0, & x \le -1.5 \\
x - \frac{x^3}{3}, & -1.5 < x < 1.5 \\
1.0, & x \ge 1.5
\end{cases}
\end{equation}

\section{Harmonic Gain Compensation \& Peak Normalization}
To maintain consistent output loudness across drive settings $G \in [1.0, 100.0]$, an dynamic automatic gain compensation multiplier $A(G)$ is applied:
\begin{equation}
A(G) = \frac{1}{\max\left(1.0, \, \tanh(G)\right)} \approx \frac{1.0}{\sqrt{G}}
\end{equation}
The final normalized output $y_{\text{final}}(n)$ is:
\begin{equation}
y_{\text{final}}(n) = A(G) \cdot y_{\text{raw}}(n)
\end{equation}

\end{document}
